% GENERATED FILE. Edit canonical lessons, then run npm run generate:books.
% canonical-source-hash: fnv1a64:f249a97fc5365619
% canonical-lessons: JA-C286-akachan, JA-C286-umareru, JA-C286-shussan, JA-C286-omedetou, JA-C286-mago

\chapter{A Baby Is Born}
\label{ch:ja-a-baby-is-born}
\hlchaptermodality{286}

\begin{chapteropening}
\textbf{By the end of this chapter:} \emph{I can say when a baby was born and wish the family well.}

Complete the last lesson of chapter 286: I can say when a baby was born and wish the family well.
\end{chapteropening}

\section[\texorpdfstring{\ja{あかちゃん} (akachan)}{あかちゃん (akachan)}]{\ja{あかちゃん} (akachan) — a baby}

\label{lesson:JA-C286-akachan}



\begin{quote}
Before the new one: say the Japanese for a commemoration, then the Japanese for a public holiday.
\end{quote}

\subsection*{You'll want to know: \ja{あかちゃん}}
\textbf{\ja{あかちゃん}} — \emph{akachan} — \textquotedblleft{}a baby\textquotedblright{}.

\begin{quote}
\textbf{\ja{あかちゃんが} \ja{ないています。}} — \emph{akachan ga naite imasu} — \textquotedblleft{}The baby is crying.\textquotedblright{}
\end{quote}

\begin{cousinweb}
\textbf{\ja{あかちゃん}} (\emph{akachan}) is a native Japanese word (wago), inherited rather than borrowed from Chinese, built from \emph{aka}, \textquotedblleft{}red\textquotedblright{} (from \emph{akai}), and \emph{-chan}, \textquotedblleft{}the affectionate suffix\textquotedblright{} (JMdict, KANJIDIC2).
\end{cousinweb}

\begin{grammarlens}[title={What you've built}]
One more word for telling what happened.
\end{grammarlens}

\subsection*{Your turn}
\begin{itemize}
\raggedright
  \item \emph{Say it:} \emph{akachan}
  \item \emph{Say it:} \emph{akachan}, once more
  \item \emph{Say it:} \emph{akachan ga naite imasu} — The baby is crying.
  \item \emph{Recall:} say the Japanese for a commemoration, then the Japanese for a public holiday, then say \emph{akachan} again
\end{itemize}

\begin{tcolorbox}[breakable,colback=teal!4,colframe=teal!35!black,title={Before you move on}]
What does \ja{あかちゃん} mean? (\textquotedblleft{}A baby\textquotedblright{}.) Say it once more.
\end{tcolorbox}
\section[\texorpdfstring{\ja{うまれる} (umareru)}{うまれる (umareru)}]{\ja{うまれる} (umareru) — to be born}

\label{lesson:JA-C286-umareru}



\begin{quote}
Before the new one: say the Japanese for a public holiday, then the Japanese for a baby.
\end{quote}

\subsection*{You'll want to know: \ja{うまれる}}
\textbf{\ja{うまれる}} — \emph{umareru} — \textquotedblleft{}to be born\textquotedblright{}.

\begin{quote}
\textbf{\ja{いもうとが} \ja{うまれた} \ja{とき、}\ja{月}\ja{が} \ja{きれいでした。}} — \emph{imōto ga umareta toki, tsuki ga kirei deshita} — \textquotedblleft{}When my little sister was born, the moon was beautiful.\textquotedblright{}
\end{quote}

Polite \textbf{\ja{うまれます}}, \emph{umaremasu}; an ichidan verb. Past (polite): \textbf{\ja{あかちゃんが} \ja{うまれました。}}, \emph{akachan ga umaremashita}.

\begin{cousinweb}
\textbf{\ja{うまれる}} (\emph{umareru}) is a native Japanese word (wago), inherited rather than borrowed from Chinese; its kanji was attached later as a kun'yomi (JMdict).
\end{cousinweb}

\begin{grammarlens}[title={What you've built}]
One more word for telling what happened.
\end{grammarlens}

\subsection*{Your turn}
\begin{itemize}
\raggedright
  \item \emph{Say it:} \emph{umareru}
  \item \emph{Say it:} \emph{umareru}, once more
  \item \emph{Say it:} \emph{imōto ga umareta toki, tsuki ga kirei deshita} — When my little sister was born, the moon was beautiful.
  \item \emph{Recall:} say the Japanese for a public holiday, then the Japanese for a baby, then say \emph{umareru} again
\end{itemize}

\begin{tcolorbox}[breakable,colback=teal!4,colframe=teal!35!black,title={Before you move on}]
What does \ja{うまれる} mean? (\textquotedblleft{}To be born\textquotedblright{}.) Say it once more.
\end{tcolorbox}
\section[\texorpdfstring{\ja{しゅっさん} (shussan)}{しゅっさん (shussan)}]{\ja{しゅっさん} (shussan) — childbirth; giving birth}

\label{lesson:JA-C286-shussan}



\begin{quote}
Before the new one: say the Japanese for a baby, then the Japanese for to be born.
\end{quote}

\subsection*{You'll want to know: \ja{しゅっさん}}
\textbf{\ja{しゅっさん}} — \emph{shussan} — \textquotedblleft{}childbirth; giving birth\textquotedblright{}.

\begin{quote}
\textbf{\ja{あねは} \ja{らいげつ} \ja{しゅっさんの} \ja{よていです。}} — \emph{ane wa raigetsu shussan no yotei desu} — \textquotedblleft{}My elder sister is due to give birth next month.\textquotedblright{}
\end{quote}

\begin{cousinweb}
\textbf{\ja{しゅっさん}} (\emph{shussan}) is Sino-Japanese (kango): two Chinese characters read with their on'yomi, readings borrowed from Chinese. They are \emph{shutsu}, \textquotedblleft{}exit\textquotedblright{} (Middle Chinese \emph{tsyhwit}) and \emph{san}, \textquotedblleft{}give birth\textquotedblright{} (\emph{sreanX}) (KANJIDIC2; Middle Chinese in Baxter's transcription). The first is shortened from \emph{shutsu} to \emph{shus-} before \emph{san}.
\end{cousinweb}

\begin{grammarlens}[title={What you've built}]
One more word for telling what happened.
\end{grammarlens}

\subsection*{Your turn}
\begin{itemize}
\raggedright
  \item \emph{Say it:} \emph{shussan}
  \item \emph{Say it:} \emph{shussan}, once more
  \item \emph{Say it:} \emph{ane wa raigetsu shussan no yotei desu} — My elder sister is due to give birth next month.
  \item \emph{Recall:} say the Japanese for a baby, then the Japanese for to be born, then say \emph{shussan} again
\end{itemize}

\begin{tcolorbox}[breakable,colback=teal!4,colframe=teal!35!black,title={Before you move on}]
What does \ja{しゅっさん} mean? (\textquotedblleft{}Childbirth; giving birth\textquotedblright{}.) Say it once more.
\end{tcolorbox}
\section[\texorpdfstring{\ja{おめでとう} (omedetō)}{おめでとう (omedetō)}]{\ja{おめでとう} (omedetō) — congratulations}

\label{lesson:JA-C286-omedetou}



\begin{quote}
Before the new one: say the Japanese for to be born, then the Japanese for childbirth.
\end{quote}

\subsection*{You'll want to know: \ja{おめでとう}}
\textbf{\ja{おめでとう}} — \emph{omedetō} — \textquotedblleft{}congratulations\textquotedblright{}.

\begin{quote}
\textbf{\ja{たんじょうび、おめでとう。}} — \emph{tanjōbi, omedetō} — \textquotedblleft{}Happy birthday!\textquotedblright{}
\end{quote}

\begin{cousinweb}
\textbf{\ja{おめでとう}} (\emph{omedetō}) is a native Japanese word (wago), inherited rather than borrowed from Chinese, and it is usually written in kana (JMdict).
\end{cousinweb}

\begin{grammarlens}[title={What you've built}]
One more word for telling what happened.
\end{grammarlens}

\subsection*{Your turn}
\begin{itemize}
\raggedright
  \item \emph{Say it:} \emph{omedetō}
  \item \emph{Say it:} \emph{omedetō}, once more
  \item \emph{Say it:} \emph{tanjōbi, omedetō} — Happy birthday!
  \item \emph{Recall:} say the Japanese for to be born, then the Japanese for childbirth, then say \emph{omedetō} again
\end{itemize}

\begin{tcolorbox}[breakable,colback=teal!4,colframe=teal!35!black,title={Before you move on}]
What does \ja{おめでとう} mean? (\textquotedblleft{}Congratulations\textquotedblright{}.) Say it once more.
\end{tcolorbox}
\section[\texorpdfstring{\ja{まご} (mago)}{まご (mago)}]{\ja{まご} (mago) — a grandchild}

\label{lesson:JA-C286-mago}



\begin{quote}
Before the new one: say the Japanese for childbirth, then the Japanese for congratulations.
\end{quote}

\subsection*{You'll want to know: \ja{まご}}
\textbf{\ja{まご}} — \emph{mago} — \textquotedblleft{}a grandchild\textquotedblright{}.

\begin{quote}
\textbf{\ja{そふと} \ja{そぼに} \ja{まごが} \ja{うまれました。}} — \emph{sofu to sobo ni mago ga umaremashita} — \textquotedblleft{}My grandparents have a new grandchild.\textquotedblright{}
\end{quote}

\begin{cousinweb}
\textbf{\ja{まご}} (\emph{mago}) is a native Japanese word (wago), inherited rather than borrowed from Chinese; its kanji was attached later as a kun'yomi (JMdict).
\end{cousinweb}

\begin{grammarlens}[title={What you've built}]
One more word for telling what happened.
\end{grammarlens}

\subsection*{Your turn}
\begin{itemize}
\raggedright
  \item \emph{Say it:} \emph{mago}
  \item \emph{Say it:} \emph{mago}, once more
  \item \emph{Say it:} \emph{sofu to sobo ni mago ga umaremashita} — My grandparents have a new grandchild.
  \item \emph{Recall:} say the Japanese for childbirth, then the Japanese for congratulations, then say \emph{mago} again
\end{itemize}

\begin{tcolorbox}[breakable,colback=teal!4,colframe=teal!35!black,title={Before you move on}]
What does \ja{まご} mean? (\textquotedblleft{}A grandchild\textquotedblright{}.) Say it once more.
\end{tcolorbox}
